% !TeX root = ../main.tex
\chapter{خوشه‌بندی و کشف ناهنجاری}
\source{اسلاید خوشه‌بندی، PDF ۳ تا ۶۴؛ کتاب فصل ۵، صفحات چاپی ۳۰۷ تا ۳۷۸؛ فصل ۹، صفحات ۷۲۳ تا ۷۵۱؛ بخش ۲٫۴. جنگل جداسازی از اسلایدها توضیح داده شده است.}
\section{هدف و نوع خوشه}

در یادگیری بدون‌نظارت برچسبِ پاسخ نداریم؛ می‌خواهیم ساختار داده را کشف کنیم. شباهت درون گروه زیاد و تمایز گروه‌ها زیاد باشد، اما این تعریف به نوع فاصله و هدف تحلیل وابسته است. شش گروه کوچک مشتری ممکن است برای بازاریابی دقیق مفید باشند و ادغامشان به دو گروه برای تصمیم مدیریتی مناسب‌تر باشد؛ وجود یک پاسخ بصری یکتا را نباید فرض کرد. کتاب بخش ۵٫۱ میان خوشه‌های جداشده، مرکزگرا، پیوسته، چگالی‌محور و مفهومی تفاوت می‌گذارد.

خوشه‌بندی افرازی یک تقسیم در یک سطح می‌سازد؛ خوشه‌بندی سلسله‌مراتبی خانواده‌ای از گروه‌های تودرتو می‌دهد که با دندروگرام نمایش داده می‌شوند. روش‌هایی مثل DBSCAN می‌توانند بعضی نقاط را نویز اعلام کنند؛ بنابراین عبارت «همهٔ نقاط حتماً عضو خوشه‌اند» ویژگی k-means است و برای همهٔ روش‌ها صدق نمی‌کند. کاربردهای اسلاید شامل بخش‌بندی مشتری، تطبیق پیشنهادهای یک محصول و تشخیص بخش‌های تصویر است. مفیدبودن خوشه‌ها باید با سؤال کاربرد سنجیده شود.

\section{k-means: الگوریتم و دلیل میانگین}

برای دادهٔ عددی، k-means مرکزهای \(\mu_1,\ldots,\mu_k\) و تقسیم را برای کم‌کردن \[
J=\sum_{j=1}^k\sum_{\boldsymbol x_i\in C_j}\|\boldsymbol x_i-\boldsymbol\mu_j\|_2^2
\] جست‌وجو می‌کند. ابتدا k مرکز تعیین می‌کنیم. در گام انتساب هر نقطه به نزدیک‌ترین مرکز می‌رود. در گام به‌روزرسانی هر مرکز میانگین اعضای جدیدش می‌شود. دو گام را تا ثابت‌شدن انتساب‌ها، جابه‌جایی ناچیز مرکزها یا رسیدن به سقف تکرار ادامه می‌دهیم. وقتی عضویت‌ها ثابت‌اند، مشتق خطای خوشه نسبت به مرکز صفر می‌شود و پاسخ میانگین است: \[
\frac{\partial}{\partial\mu_j}\sum_{x_i\in C_j}(x_i-\mu_j)^2=2|C_j|\mu_j-2\sum_{x_i\in C_j}x_i=0.
\]

در مثال یک‌بعدی \(1,2,8,9\)، $k=2$ و مرکزهای اولیهٔ ۱ و ۸، انتساب‌ها \(\{1,2\}\) و \(\{8,9\}\) می‌شوند. مرکزهای جدید ۱٫۵ و ۸٫۵ و SSE برابر ۱ است. برای $k=1$ مرکز ۵ و SSE برابر ۵۰ است؛ برای $k=4$، SSE صفر می‌شود. بنابراین کاهش خام SSE با افزایش k دلیل کافی برای انتخاب تعداد بیشتر خوشه نیست.

هر گام، با قراردادهای ثابت، هدف را افزایش نمی‌دهد و الگوریتم به یک جواب محلی می‌رسد؛ تضمین بهترین جواب جهانی وجود ندارد. بذرهای متفاوت می‌توانند پاسخ‌های متفاوت بدهند. چند شروع مستقل و انتخاب کمترین SSE برای k ثابت مفید است. کتاب در چاپی ۳۲۵--۳۲۶/PDF ۳۲۶--۳۲۷ روش k-means++ را بسط می‌دهد: اولین مرکز تصادفی است و مرکزهای بعدی با احتمال متناسب با مربع فاصله از نزدیک‌ترین مرکز انتخاب می‌شوند. خوشهٔ خالی را باید با قرارداد مشخص، مثلاً قراردادن مرکز روی نقطهٔ دارای بیشترین فاصله، مدیریت کرد.

زمان تقریبی \(O(nkpI)\) برای n نقطه، p ویژگی و I تکرار است. استانداردسازی ویژگی‌ها اهمیت دارد. حساسیت به نقاط دور، شکل‌های غیرکروی و گروه‌های با اندازه یا چگالی بسیار متفاوت از محدودیت‌های مهم‌اند. منحنی آرنج محل کاهش منفعت افزودن k را نشان می‌دهد ولی همیشه آرنج آشکاری وجود ندارد. اسلاید ۲۱ X-means را به‌عنوان گسترش انتخاب تعداد خوشه معرفی می‌کند؛ این روش با انتخاب سادهٔ آرنج یکسان نیست.

\section{k-medoids و تمایز میانه از medoid}

اسلاید ۲۶ k-medoids را با عبارت «استفاده از میانه» توضیح می‌دهد، اما تعریف دقیق این است که نماینده باید یک مشاهدهٔ موجود باشد: \[
\operatorname{medoid}(C)=\arg\min_{z\in C}\sum_{x\in C}d(x,z).
\] میانهٔ هر مؤلفه در فضای چندبعدی الزاماً نقطه‌ای از داده نیست و k-medians با k-medoids متفاوت است. در یک‌بعد با فاصلهٔ مطلق، یکی از میانه‌های مشاهده‌شده می‌تواند medoid باشد. مثال \(1,3,5,7,1009\) میانگین ۲۰۵ ولی میانه و medoid برابر ۵ می‌دهد. مزیت نسبت به میانگینِ مربع‌خطا، مقاومت بیشتر در برابر نقاط بسیار دور است؛ ولی انتخاب نماینده می‌تواند هزینهٔ محاسباتی بیشتری داشته باشد.

\section{silhouette: سنجش انتساب هر نقطه}

برای نقطه i ، \(a_i\) میانگین فاصله تا سایر اعضای خوشهٔ خودش و \(b_i\) کمترین میانگین فاصله تا اعضای یک خوشهٔ دیگر است: \[
s_i=\frac{b_i-a_i}{\max(a_i,b_i)},\qquad S=\frac1n\sum_i s_i.
\] اگر \(a=2,b=6\)، مقدار \(2/3\) نشانهٔ انتساب مناسب است؛ اگر \(a=6,b=2\)، مقدار \(-2/3\) نشان می‌دهد نقطه به خوشهٔ دیگر نزدیک‌تر است. نزدیک صفر یعنی مرز مبهم. برای خوشهٔ تک‌عضوی معمولاً $s=0$ قرارداد می‌شود؛ تعریف به دست‌کم دو خوشه و کمتر از n خوشه نیاز دارد. میانگین بالا مفید است اما «مستقل از تعداد خوشه» را نباید به معنی بی‌تأثیربودن k دانست: خود نمره با تقسیم‌بندی، مقیاس و شکل داده عوض می‌شود و گروه‌های کروی و جدا را ترجیح می‌دهد. بهتر است نمودار silhouette هر خوشه نیز کنار میانگین دیده شود.

\section{DBSCAN و اتصال نواحی پرتراکم}

اگر شکل داده هلالی باشد، مرکز کروی k-means ابزار مناسبی نیست. DBSCAN خوشه را ناحیهٔ پرتراکمِ متصل می‌بیند. همسایگی \(N_\varepsilon(x)=\{z:d(x,z)\le\varepsilon\}\) شامل خود نقطه است. نقطهٔ هسته حداقل MinPts عضو در این همسایگی دارد؛ نقطهٔ مرزی هسته نیست ولی همسایهٔ یک هسته است؛ بقیه نویزند. مرز فقط به خوشه متصل می‌شود و خودش اجازهٔ گسترش اتصال را نمی‌دهد.

روند آموزشی روشن: ابتدا هسته‌ها را تعیین کنیم، میان هسته‌هایی که فاصله‌شان حداکثر $\varepsilon$ است یال بگذاریم، مؤلفه‌های متصلِ هسته‌ها را خوشه بگیریم و سپس مرزی‌ها را وصل کنیم. اگر مرزی به دو مؤلفه نزدیک باشد، قرارداد ترتیب/شکستن تساوی ممکن است بر عضویتش اثر بگذارد؛ مؤلفه‌های هسته به همین دلیل نباید از راه یک مرزی ادغام شوند. این همان بسط بخش ۵٫۴٫۲ کتاب است.

مثال \(x=(0,0.1,0.2,0.35,2)\)، $\varepsilon=0.16$ و $\mathrm{MinPts}=3$: نقطهٔ ۰٫۱ همسایه‌های ۰،۰٫۱،۰٫۲ و نقطهٔ ۰٫۲ همسایه‌های ۰٫۱،۰٫۲،۰٫۳۵ را دارند و هسته‌اند. ۰ و ۰٫۳۵ مرزی و ۲ نویز است. در نتیجه یک خوشهٔ چهارعضوی می‌سازیم. زیادکردن $\varepsilon$ ممکن است خوشه‌های جدا را وصل کند؛ کم‌کردن آن یا زیادکردن MinPts ممکن است نویز را زیاد کند. نمودار مرتب فاصله تا همسایهٔ k ام راهنمای $\varepsilon$ است، اما قرارداد شمردن خود نقطه باید شفاف باشد. چگالی‌های بسیار متفاوت و ابعاد بالا نقاط ضعف اصلی‌اند. پیچیدگیِ جست‌وجوی ساده مربعی است؛ نمایهٔ مکانی در ابعاد کم می‌تواند کمک کند.

\section{سلسله‌مراتبی: محتوای مکملِ تصریح‌شده در درس}

اسلاید ۶۴ این مبحث را به منبع مکمل ارجاع می‌دهد؛ توضیح زیر بر بخش ۵٫۳ کتاب استوار است. در روش تجمعی هر نقطه ابتدا یک خوشه است؛ نزدیک‌ترین دو خوشه ادغام و فاصله‌ها به‌روز می‌شوند. ارتفاع ادغام در دندروگرام، مقدار ناهمانندیِ اتصال است؛ برش افقی درخت یک تقسیم می‌دهد.

single linkage کمترین فاصلهٔ دو عضو، complete بیشترین فاصله و average میانگین همهٔ فاصله‌های بین‌خوشه‌ای است. Ward ادغامی را برمی‌گزیند که کمترین افزایش SSE را دارد. برای دو مجموعهٔ یک‌بعدی $A=\{0,2\}$ و $B=\{5,9\}$، فاصله‌های جفتی ۵،۹،۳،۷ است؛ $d_{\mathrm{single}}=3$، $d_{\mathrm{complete}}=9$، $d_{\mathrm{average}}=6$. پس انتخاب اتصال واقعاً ساختار خروجی را عوض می‌کند. Single می‌تواند گروه‌های کشیده را پیدا کند ولی به زنجیرهٔ نقاط نویزی حساس است؛ complete گروه‌های فشرده‌تر می‌سازد. ادغام‌های حریصانه برگشت‌ناپذیرند و حافظهٔ ماتریس فاصله می‌تواند \(O(n^2)\) باشد.

\section{فاصلهٔ مناسب، بخشی از تعریف مسئله است}

برای p ویژگی عددی، فاصلهٔ مینکوفسکی \[
d_q(x,z)=\left(\sum_{j=1}^p|x_j-z_j|^q\right)^{1/q},\quad q\ge1
\] است؛ $q=1$ منهتن و $q=2$ اقلیدسی. برای اختلاف مؤلفه‌ها ۳ و ۴، این دو فاصله به‌ترتیب ۷ و ۵ هستند. ریشه در تعریف فاصلهٔ اقلیدسی ضروری است. برای صفت اسمی، برابری/نابرابری طبیعی‌تر از تفاضل عددِ کدگذاری است؛ برای ترتیبی می‌توان رتبه‌ها را با قرارداد نرمال کرد. فاصلهٔ لونشتاین حداقل درج/حذف/جایگزینی حرف است؛ table و cable فاصلهٔ ۱ دارند. اسلاید ۴۰ دیگر خانواده‌های رشته‌ای را فقط معرفی می‌کند، پس تبدیل آن به فصل عمیق پردازش زبان خارج از دامنهٔ فعلی است.

به‌عنوان بسط کتاب، Jaccard برای مجموعه‌ها \(|A\cap B|/|A\cup B|\) و cosine برای بردارهای غیرصفر \(x^Tz/(\|x\|\|z\|)\) است. نبود مشترکِ یک رخداد نادر معمولاً در Jaccard پاداش نمی‌گیرد. انتخاب فاصله باید با مفهوم شباهت در کاربرد هماهنگ باشد؛ استانداردسازی به‌خودی‌خود فاصلهٔ نامناسب را مناسب نمی‌کند.

\section{مسئله و روش‌های مقاوم}

یک مشاهده وقتی ناهنجار است که با ساختار معمول مورد انتظار ناسازگار باشد؛ این ناسازگاری می‌تواند زمینه‌ای یا محلی باشد. دمای ۳۰ درجه در تابستان طبیعی و برای یک سردخانه غیرعادی است. دادهٔ نادر الزاماً خطا نیست: ممکن است همان تقلب یا رخداد مورد نظر باشد. پس حذف همهٔ برون‌افتاده‌ها سیاست عمومی نیست. اسلاید ۴۴ تأکید دارد اگر پالایش آموزش انتخاب شد، آزمون برای ارزیابی واقعی دست‌نخورده بماند؛ آستانه‌ها نیز از آموزش یاد گرفته شوند.

حصار توکی با \(IQR=Q_3-Q_1\) و بازهٔ \([Q_1-1.5IQR,Q_3+1.5IQR]\) یک قاعدهٔ اکتشافی مقاوم است و برخلاف عبارت اسلاید ۴۶، نرمال‌بودن پیش‌شرط محاسبهٔ آن نیست. انتخاب روش محاسبهٔ چارک، به‌ویژه در نمونهٔ کوچک، جواب را تغییر می‌دهد. برای دادهٔ مرتب \((0,1,2,3,4,5,6,20)\)، با میانهٔ نیمه‌ها $Q_1=1.5$ و $Q_3=5.5$، پس $IQR=4$ و حصارها −4.5 و 11.5 هستند؛ ۲۰ علامت می‌خورد. این مثال قرارداد را صریح دارد و از ابهام چارک اسلاید جلوگیری می‌کند.

انحراف مطلق میانه \(MAD=\operatorname{median}_i|x_i-\operatorname{median}(X)|\) است. در مثال اسلاید \((0,1,1,3,5,7,42)\)، میانه ۳ و $MAD=2$ است؛ با قاعدهٔ خام $k=3$ بازهٔ −۳ تا ۹ می‌شود و ۴۲ بیرون است. MAD خام با تخمین سازگار انحراف معیار نرمال \(1.4826\,MAD\) فرق دارد؛ آستانه‌ها را نباید مخلوط کرد. اگر $MAD=0$ شد تقسیم بر آن معتبر نیست و باید قرارداد مقاوم جایگزین یا گزارش عدم امکان تعیین نمره داشت. IQR/MAD تک‌متغیره ممکن است نقطه‌ای در «حفرهٔ میان دو گروه» را نبینند؛ مقدارِ عددی آن نقطه شدید نیست ولی از همهٔ همسایگان دور است.

\section{روش آماری، فاصله‌ای و چگالی نسبی}

در مدل نرمال تک‌متغیره، نمرهٔ استاندارد \(z=(x-\mu)/\sigma\) دوری از مرکز را نسبت به پراکندگی می‌سنجد. ارزیابی احتمال به درست‌بودن مدل توزیع بستگی دارد؛ برای دادهٔ چندقله‌ای یک نرمال واحد ممکن است گمراه کند. در چندبعد با ویژگی‌های هم‌بسته، فاصلهٔ ماهالانوبیسِ مربعی \((x-\mu)^T\Sigma^{-1}(x-\mu)\) شکل و جهت ابر داده را لحاظ می‌کند؛ کوواریانس باید قابل‌برآورد و غیرتکین یا منظم‌شده باشد.

روش‌های فاصله‌ای اسلاید ۴۹ سه صورت دارند: تعداد همسایه در شعاع ثابت، فاصله تا همسایهٔ k ام، یا میانگین فاصله تا k همسایه. فاصلهٔ بیشتر می‌تواند نمرهٔ ناهنجاری بزرگ‌تر بدهد. اما یک خوشهٔ طبیعی کم‌تراکم نباید صرفاً به‌علت فاصله‌های بزرگ‌تر غیرعادی شود. اینجا مقایسهٔ محلی کمک می‌کند.

کتاب PDF ۷۴۳ نسبت «میانگین چگالی همسایه‌ها / چگالی نقطه» را برای شهود توضیح می‌دهد و صریحاً می‌گوید LOF دقیق از تعریف مقاوم‌تری استفاده می‌کند. برای معرفی دقیق LOF باید نوشت: \[
\operatorname{reachdist}_k(x,o)=\max\{k\text{-dist}(o),d(x,o)\}.
\] \[
\operatorname{lrd}_k(x)=\left(\frac{\sum_{o\in N_k(x)}\operatorname{reachdist}_k(x,o)}{|N_k(x)|}\right)^{-1},
\] \[
\operatorname{LOF}_k(x)=\frac{1}{|N_k(x)|}\sum_{o\in N_k(x)}\frac{\operatorname{lrd}_k(o)}{\operatorname{lrd}_k(x)}.
\] نزدیک ۱ یعنی تراکم نقطه با همسایه‌ها سازگار است؛ مقدار خیلی بزرگ‌تر نامعمول است، ولی هر مقدار صرفاً بزرگ‌تر از ۱ حکم قطعی ناهنجاری نیست. انتخاب k و آستانه به مسئله وابسته است. پیاده‌سازی آزمایشگاه اگر فقط عکس فاصلهٔ میانگین را استفاده کند باید «چگالی نسبی ساده» نام بگیرد، نه LOF دقیق.

\section{خوشه‌محور و جنگل جداسازی}

DBSCAN مستقیماً برچسب نویز می‌دهد، اما نویز با پارامترهای انتخاب‌شده تعریف شده و اثبات خطای داده نیست. روش خوشه‌محور می‌تواند فاصله از مرکز، فاصلهٔ نسبی به پراکندگی همان خوشه، کوچک‌بودن گروه و دوری آن از گروه‌های بزرگ را ترکیب کند. اسلایدهای ۵۳--۵۵ CBLOF را با اندازهٔ خوشه و فاصله تا خوشهٔ بزرگ معرفی می‌کنند. کتاب بخش ۹٫۵ بحث متناظر خوشه‌ای دارد، اما فرمول دقیق CBLOF را ارائه نمی‌کند؛ فرمول ساختگی به آن نسبت داده نشود. گروه کوچک می‌تواند یک جمعیت واقعی باشد و باید زمینهٔ کاربرد را دید.

Isolation Forest به‌جای مدل‌کردن چگالی، می‌پرسد چند برش تصادفی برای جداکردن نقطه لازم است. در هر گره یک ویژگی و سپس یک مقدار شکستن بین کمینه و بیشینهٔ همان ویژگی تصادفی انتخاب می‌شود. این روند روی زیرگروه‌ها تکرار می‌شود و عمقِ مسیر نقطه ثبت می‌شود. نقطهٔ دورافتاده اغلب با برش‌های کمتری تنها می‌شود؛ نقاط درون ناحیهٔ انبوه به برش‌های بیشتر نیاز دارند. میانگین‌گیری طول مسیر در چند درخت، تصادف یک درخت را تعدیل می‌کند. نسخهٔ آموزشی اسلاید ۵۶--۶۳ همین شهود و احتمال‌های جداسازیِ یک‌گامه را نشان می‌دهد؛ طول مسیر کوتاه‌تر یعنی ناهنجاری بیشتر. برای توضیح این مفهوم نیازی به معرفی نمرهٔ نرمال‌شده‌ای که در منابع ارسالی ارائه نشده نیست. سقف عمق و نمونه‌گیری، الگوریتم عملی را سریع‌تر می‌کنند.
\section{اجرای دستی؛ چه چیزی در هر گام تغییر می‌کند؟}
برای نقاط یک‌بعدی $1,2,8,9$، دو مرکز اولیهٔ ۱ و ۸ انتخاب کنید. نزدیک‌ترین مرکز، دو گروه $\{1,2\}$ و $\{8,9\}$ می‌دهد. پیش از به‌روزرسانی، مجموع مربع فاصله‌ها $0+1+0+1=2$ است. با ثابت ماندن عضویت‌ها، مرکزها به $1.5,8.5$ می‌روند و هدف به $4\times0.25=1$ کاهش می‌یابد. انتساب بعدی همان گروه‌ها را حفظ می‌کند؛ این اجرا همگراست. «یک گام» را در گزارش دقیق تعریف کنید: انتساب، به‌روزرسانی یا یک چرخهٔ کامل؟

در مقابل، DBSCAN مرکز و میانگین نمی‌سازد. برای نقاط $0,0.1,0.2,0.35,2$ و $\varepsilon=0.16$، حداقل همسایهٔ سه با احتساب خود نقطه، فقط $0.1$ و $0.2$ را هسته می‌کند. این دو به هم متصل‌اند؛ صفر و $0.35$ مرزی و ۲ نویز است. یک نقطهٔ مرزی اگر نزدیک دو مؤلفهٔ هسته باشد، طبق قرارداد به یکی وصل می‌شود؛ نباید دو مؤلفه را از راه آن ادغام کنیم. این تفاوت دلیل وجود دو مرحلهٔ «اتصال هسته‌ها» و «انتساب مرزی‌ها» است.

\subsection*{شکل نمودار را با تعریف فاصله بخوانید}
در نمودار اقلیدسی، یک واحد افقی و عمودی باید طول تصویری یکسان داشته باشند تا دایرهٔ همسایگی و عمود بودن تصویر PCA درست دیده شود. برای سن و درآمد، محورها واحدهای متفاوت دارند و شکل خام به‌تنهایی نمایش فاصلهٔ استانداردشده نیست. عدد فاصله را با قرارداد تبدیل بخوانید. در ناهنجاری هم باید روشن باشد که رنگ، آستانهٔ واقعی است یا فقط رتبهٔ چند نمرهٔ بزرگ‌تر؛ رتبهٔ اول بودن لزوماً به معنی خطای داده نیست.

\section{خودآزمایی و تمرین}
\begin{selftests}
\begin{selftest}{یک چرخهٔ k-means را کامل کنید}
برای $1,2,8,9$ و مرکزهای اولیهٔ ۱ و ۸، مرکزهای تازه و SSE چیست؟ آیا SSE صفر با چهار مرکز، چهار خوشه را توجیه می‌کند؟
\answer{مرکزها $1.5,8.5$ و SSE یک است. چهار مرکز می‌توانند SSE را صفر کنند، اما این فقط کاهش هدف آموزش است و تعداد گروه‌های معنادار را ثابت نمی‌کند.}
\end{selftest}
\begin{selftest}{میانه و medoid یکسان‌اند؟}
آیا میانهٔ مؤلفه‌ای در دوبعد حتماً یکی از نقاط مشاهده‌شده است؟
\answer{خیر؛ medoid باید عضو واقعی داده باشد و مجموع ناهمانندی را کمینه کند، ولی میانهٔ مؤلفه‌ای ممکن است مشاهده نشده باشد. روش‌های مبتنی بر میانه و medoid را نباید با یک نام معرفی کرد.}
\end{selftest}
\begin{selftest}{نقطهٔ مرزی را به‌جای هسته استفاده نکنید}
چرا در DBSCAN، اتصال دو خوشه از راه یک نقطهٔ غیرهسته درست نیست؟
\answer{گسترش چگالی از هسته انجام می‌شود. نقطهٔ مرزی به هسته دسترسی دارد ولی خودش تراکم کافی برای ادامهٔ زنجیره ندارد؛ استفاده از آن برای ادغام می‌تواند دو ناحیهٔ هستهٔ جدا را نادرست یکی کند.}
\end{selftest}
\begin{selftest}{معیار اتصال چه تفاوتی می‌سازد؟}
برای دو گروه $A=\{0,2\}$ و $B=\{5,9\}$، فاصلهٔ single، complete و average چیست؟
\answer{فاصله‌های جفتی $5,9,3,7$ هستند؛ کمینه ۳، بیشینه ۹ و میانگین ۶ است. انتخاب اتصال بر ارتفاع ادغام و گروه‌های نهایی اثر واقعی دارد.}
\end{selftest}
\begin{selftest}{آیا IQR نرمال بودن می‌خواهد؟}
برای دادهٔ دو‌قله‌ای، آیا محاسبهٔ حصار IQR ممنوع است؟ آیا هر نقطهٔ داخل حصار عادی است؟
\answer{محاسبهٔ چارک و حصار به نرمال بودن نیاز ندارد؛ ولی تفسیر حصار به شکل داده و زمینه وابسته است. یک نقطه در حفرهٔ میان گروه‌ها ممکن است داخل حصار تک‌متغیره باشد ولی از همسایگان دور باشد.}
\end{selftest}
\begin{selftest}{نمرهٔ محلی را درست تفسیر کنید}
آیا هر $LOF>1$ یا هر نقطهٔ با رتبهٔ بالاتر باید حذف شود؟
\answer{خیر؛ LOF نزدیکی چگالی را با همسایگان مقایسه می‌کند و آستانهٔ کاربردی جدا می‌خواهد. رتبهٔ بالا، خطای ثبت یا بی‌ارزش بودن داده را اثبات نمی‌کند؛ ممکن است رخداد مهم یا جمعیتی واقعی باشد.}
\end{selftest}
\end{selftests}
